Descrição da Atividade

Algoritmos de ordenação têm como objetivo colocar elementos
em uma sequência ordenada. Existem pelo menos quatro formas
bastante difundidas de ordenação: bubble sort, insertion sort, merge
sort e quick sort. Dependendo dos dados a serem ordenados, é
possível que um algoritmo diferente seja o mais eficiente para cada
cenário. Com base no padrão Strategy, implemente um sistema
com uma classe cliente (contexto) que pode alternar livremente a
estratégia de ordenação usada para ordenar uma lista de inteiros a
partir da seguinte assinatura de método:

List<Integer> ordena(List<Integer> elementos);

A lista recebida pelo método não pode ser modificada, e a or-
denação deve operar sempre sobre uma cópia defensiva. Cada es-
tratégia deve definir um limite conceitual de operação baseado em
um tamanho máximo de lista que impeça seu uso. Caso o limite
estabelecido pela estratégia selecionada seja violado e ela recuse o
processamento, o contexto deve automaticamente recorrer ao algo-
ritmo padrão (merge sort) como estratégia de fallback e registrar,
por meio de um print, que houve substituição da estratégia original.
Não é necessário implementar os algoritmos de ordenação.
